Les 10 AXIOMES 
---

**Axiome 1 : Accapareur urbain**
Un citoyen est considéré comme un accapareur urbain lorsqu'il possède au moins quatre parcelles urbaines.
**DL :**
AccapareurUrbain ≡ Citoyen ⊓ (≥4 possede.ParcelleUrbaine)

---

**Axiome 2 : Conflit d'intérêt direct**
Un agent public est en situation de conflit d'intérêt lorsqu'il traite un dossier dont il est lui-même bénéficiaire.
**DL :**
ConflitInteretDirect ≡ AgentPublic ⊓ ∃traite.Dossier ⊓ ∃beneficiaire.Dossier

---

**Axiome 3 : Spéculateur foncier**
Une personne est considérée comme spéculateur foncier lorsqu'elle revend une parcelle moins d'un an après son acquisition en réalisant une plus-value supérieure à 50 %.
**DL :**
SpeculateurFoncier ≡ Personne ⊓ ∃achete.Parcelle ⊓ ∃vend.Parcelle
(*les contraintes de temps et de pourcentage sont gérées hors DL*)

---

**Axiome 4 : Prête-nom**
Un citoyen est considéré comme un prête-nom lorsqu'il possède une parcelle et partage le même numéro de téléphone qu'une autre personne possédant également cette parcelle.
**DL :**
PreteNom ≡ Citoyen ⊓ ∃possede.Parcelle ⊓ ∃partageTelephone.Citoyen

---

**Axiome 5 : Réseau de corruption**
Un agent public appartient à un réseau de corruption lorsqu'il traite un dossier dont le bénéficiaire entretient un lien financier avec lui.
**DL :**
ReseauCorruption ≡ AgentPublic ⊓ ∃traite.Dossier ⊓ ∃lienFinancier.Personne

---

**Axiome 6 : Promoteur fantôme**
Un promoteur est qualifié de promoteur fantôme lorsqu'il effectue des transactions foncières suspectes sans posséder officiellement de parcelles.
**DL :**
PromoteurFantome ≡ Promoteur ⊓ ¬∃possede.Parcelle ⊓ ∃transactionSuspecte.Transaction

---

**Axiome 7 : Dossier suspect**
Un dossier actif est considéré comme suspect lorsqu'il comporte au moins trois alertes de fraude ou d'irrégularité.
**DL :**
DossierSuspect ≡ DossierActif ⊓ (≥3 contientAlerte.Alerte)

---

**Axiome 8 : Accaparement familial**
Une famille est en situation d'accaparement foncier lorsque ses membres possèdent collectivement au moins huit parcelles urbaines.
**DL :**
AccaparementFamilial ≡ Famille ⊓ (≥8 possede.ParcelleUrbaine)

---

**Axiome 9 : Blanchiment circulaire**
Il y a blanchiment circulaire lorsqu'une même parcelle est vendue successivement entre plusieurs personnes puis revient à son propriétaire initial dans un court délai.
**DL :**
BlanchimentCirculaire ≡ Parcelle ⊓ ∃cycleVente.Transaction

---

**Axiome 10 : Notaire complice**
Un notaire est considéré comme complice lorsqu'il intervient dans au moins cinq dossiers reconnus comme suspects.
**DL :**
NotaireComplice ≡ Notaire ⊓ (≥5 traite.DossierSuspect)


LES 8 CONTRAINTES D'INTEGRITES



---

### **CI-1 : Auto-traitement interdit**

Un agent public ne peut pas traiter un dossier dont il est le bénéficiaire.
**DL / règle :**
AgentPublic ⊓ ∃traite.Dossier ⊓ ∃beneficiaire.Dossier ⊑ ⊥
*(ou contrainte d’incohérence)*

---

### **CI-2 : Limite de possession de parcelles urbaines**

Un citoyen ne peut pas posséder plus de 3 parcelles urbaines.
**DL :**
Citoyen ⊓ (≥4 possede.ParcelleUrbaine) ⊑ ⊥

---

### **CI-3 : Unicité de propriété d’une parcelle**

Une parcelle ne peut pas avoir plusieurs propriétaires en même temps.
**DL :**
Parcelle ⊓ ∃possede⁻.Personne ⊓ (≥2 possede⁻.Personne) ⊑ ⊥
*(ou : fonctionnalité du rôle “possede” inversé)*

---

### **CI-4 : Conflit notaire–bénéficiaire interdit**

Un notaire ne peut pas être bénéficiaire d’un dossier qu’il traite.
**DL :**
Notaire ⊓ ∃traite.Dossier ⊓ ∃beneficiaire.Dossier ⊑ ⊥

---

### **CI-5 : Partage de téléphone suspect**

Deux personnes différentes ne peuvent pas partager un numéro de téléphone si elles possèdent une même parcelle (suspicion de prête-nom).
**DL :**
(Citoyen ⊓ ∃partageTelephone.Citoyen ⊓ ∃possede.Parcelle ⊓ ∃possede.Parcelle) ⊑ ⊥

---

### **CI-6 : Cohérence des dates de transaction**

Une vente ne peut pas avoir une date antérieure ou égale à la date d’achat.
**DL :**
Transaction ⊓ (dateVente ≤ dateAchat) ⊑ ⊥

---

### **CI-7 : Promoteur doit posséder au moins une parcelle**

Un promoteur sans parcelle est incohérent dans le système.
**DL :**
Promoteur ⊓ ¬∃possede.Parcelle ⊑ ⊥

---

### **CI-8 : Dossier suspect ne peut pas être actif**

Un dossier suspect ne peut pas rester actif dans le système.
**DL :**
DossierSuspect ⊓ DossierActif ⊑ ⊥

---

